\section{Diseño y desarrollo del sistema}
En esta etapa se transformaron los criterios teóricos y metodológicos en un sistema funcional, integrando de manera práctica el uso de Git y GitHub, la Guía de Scrum 2020, los patrones de diseño, el patrón Modelo–Vista–Controlador (MVC) y el lenguaje Java como ejes de construcción y organización del proyecto.

\subsection{Procedimiento detallado}
\subsubsection{Diseño de la arquitectura (MVC)}
Se definieron las capas principales del sistema: Modelo, encargado de la gestión de datos y la lógica de negocio; Vista, responsable de las interfaces de usuario; y Controlador, que coordina la comunicación entre ambos. Esta separación de responsabilidades permitió aumentar la mantenibilidad y la escalabilidad, facilitando la evolución independiente de cada componente.

\subsubsection{Patrones de Diseño}
Se aplicaron patrones creacionales y estructurales para optimizar la reutilización de código y garantizar una arquitectura clara y flexible. Esta práctica redujo la duplicación y facilitó la incorporación de nuevas funcionalidades.

\subsubsection{Lenguaje Java}
El desarrollo principal se realizó en Java, aprovechando su orientación a objetos, su portabilidad mediante la JVM y su amplio ecosistema de librerías. Se emplearon principios de encapsulamiento, herencia y polimorfismo para asegurar un código robusto y extensible.

\subsubsection{Scrum}
El equipo trabajó en sprints definidos, con reuniones de planificación, revisiones y retrospectivas. La auto-gestión y la transparencia recomendada por la guía 2020 permitieron ajustar prioridades y detectar incidentes de forma temprana, fomentando la mejora continua.

\subsubsection{Git and Github}
Todo el código se gestionó con Git y se alojó en GitHub, lo que brindó trazabilidad, trabajo en ramas, revisiones mediante pull requests y una integración continua eficaz para cada entrega.

\subsection{Herramientas complementarias}
• Backend: Java con Spring Boot, elegido por su seguridad y escalabilidad.

• Frontend web: React.js, por su flexibilidad en interfaces dinámicas.

• Aplicación móvil: React Native, para extender funcionalidades a entornos móviles sin duplicar código.

• Base de datos: PostgreSQL, encargada del manejo de la información de usuarios, eventos y reservas.

Esta fase consolidó la integración de herramientas, metodologías y arquitecturas, permitiendo obtener un sistema moderno, escalable y alineado con las mejores prácticas del desarrollo de software.